\section{Estimating target composition protection from the observed data}
\label{sec:target-protection}

The universal CATE range 2 can dominate population interval width. We now
construct one target-outcome confidence event that simultaneously controls the
range and empirical variance of the true CATE. The target-only comparator uses
the identical event and composition method. No simulated true range or slope
enters the procedure.

Let $b=(\beta_t^1-\beta_t^0)_{1:p}$ denote the slope contrast, excluding the
intercept. The target features have a \emph{predeclared} rectangular support
$\prod_{r=1}^p[l_r,u_r]$; observed minima and maxima do not define that support.
Let $X_{ta}$ include its intercept, and let $\widehat b$ be the difference of
its two OLS slope vectors. Conditional on all target designs,
\[
 V_b=\tfrac14\left\{(X_{t1}^TX_{t1})^{-1}+
                         (X_{t0}^TX_{t0})^{-1}\right\}_{1:p,1:p}
\]
is a sub-Gaussian variance proxy, since outcomes lie in $[0,1]$ and the target
conditional means are linear. Thus, for every vector $v$,
\[
 \E\{\exp[v^T(\widehat b-b)]\mid\mathcal D\}
 \le\exp(\tfrac12v^TV_bv).
\]
The proxy is determined by the design, not estimated from outcome residuals.

Put $x_R=\log(1/\alpha_R)$ and
$r_R^2=p+2\sqrt{p x_R}+2x_R$. A sub-Gaussian quadratic-form bound
\citep{hsu2012} gives the conditional confidence region
\[
 \mathcal C_b=\{v:(v-\widehat b)^TV_b^{-1}(v-\widehat b)\le r_R^2\},
 \qquad\Prb(b\in\mathcal C_b\mid\mathcal D)\ge1-\alpha_R.
\]
This is a sub-Gaussian confidence ellipsoid, not an exact Gaussian confidence
region. Applying the bound to $V_b^{-1/2}(\widehat b-b)$ proves the claim.

\subsection{One event for both range and variance}

Let $w_r=u_r-l_r$. On the same confidence event, the CATE range is at most
\[
 R^+=\min\left\{2,\max_{s\in\{-1,1\}^p}
   \left[(w\odot s)^T\widehat b+
        r_R\sqrt{(w\odot s)^TV_b(w\odot s)}\right]\right\}.
\]
This follows by maximizing the rectangular-support range
$\sum_r w_r|b_r|$ over the ellipsoid. The implementation enumerates signs through
$p=12$; above that it uses the valid upper bound
$\sum_r w_r(|\widehat b_r|+r_R\sqrt{V_{b,rr}})$, capped at 2.
The first experiments focus on $p=4,10$.

Let $S_X$ be the target sample covariance of features with denominator $n_t-1$.
Writing $V_b=L_bL_b^T$, the true empirical CATE variance is bounded by
\[
 V_X^+=\min\left\{
  \left[\sqrt{\widehat b^TS_X\widehat b}
       +r_R\sqrt{\lambda_{\max}(L_b^TS_XL_b)}\right]^2,
  \frac{n_t}{n_t-1}\frac{(R^+)^2}{4}\right\}.
\]
The first term is the triangle inequality in the $S_X$ seminorm; the second is
the bounded-range sample-variance inequality. The observed matrix $S_X$ does
not require a separate confidence event: both statements hold deterministically
for every $b\in\mathcal C_b$.

\subsection{Population enlargement and error allocation}

For the true, fixed CATE function, the empirical Bernstein inequality of
\citet[Theorem 4]{maurer2009}, applied in both directions, gives
\[
 r_X^{\rm EB}=\sqrt{\frac{2V_X^+\log(4/\alpha_X)}{n_t}}
        +\frac{7R^+\log(4/\alpha_X)}{3(n_t-1)}.
\]
The range-only alternative is
$r_X^{\rm H}=R^+\sqrt{\log(2/\alpha_X)/(2n_t)}$.
Although their feasible bounds depend on target outcomes, coverage follows
from the union of the coefficient event and the concentration event for the
\emph{true} CATE. It does not require independence between those events.
An adaptive option takes their minimum only after allocating $\alpha_X/2$ to
each inequality. Taking the unadjusted minimum is not justified.

For the source-assisted procedure, the fixed allocation is
\[
 \alpha_R=.005,\quad\alpha_X=.01,\quad
 \alpha_S=.01,\quad\alpha_T=.01,\quad\alpha_D=.015.
\]
The new universal-range, range-only, empirical-Bernstein and adaptive comparisons
share exactly the same conditional source and target bands. The coefficient
allowance is unused in the universal-range control. The v19 allocation is
retained separately as a historical comparator, rather than silently changed.
For standalone target-only, use the same $.005$ coefficient event and $.0225$
each for outcomes and composition. This preserves a common $.05$ total failure
budget while improving the comparator together with the proposed method.

These certificates address target composition, not arbitrary target outcome
misspecification. At larger $p$, the coefficient region can be too wide for a
useful range bound. Variance-sensitive composition can still improve before
range estimation becomes precise; the experiments report both components.

\subsection{Adaptation when the target effect is nearly constant}

For fixed $p$, bounded support and well-conditioned target designs, the proxy
$V_b$ is $O(n_t^{-1})$. The coefficient region consequently contributes
$O_p(n_t^{-1/2})$ to the empirical CATE standard-deviation bound. If
$\sigma_\tau^2=\Var_t\{\tau(X)\}$, this gives
\[
 r_X^{\rm EB}=O_p(\sigma_\tau n_t^{-1/2}+n_t^{-1}).
\]
In particular, a constant CATE has $b=0$, $R^+=O_p(n_t^{-1/2})$ and
$V_X^+=O_p(n_t^{-1})$, so the composition allowance is $O_p(n_t^{-1})$.
The algorithm is not told that the effect is constant.

The same order in expectation on good designs follows from conditional
second-moment bounds for the coefficient error and Jensen's inequality for
its $S_X$ seminorm. For the true CATE, the sample variance is unbiased and its
expected square root is at most $\sigma_\tau$. Combine the good-design
expectation with the bounded-width fallback term in
Section~\ref{sec:design-stability}. These statements require the stated
fixed-dimension/conditioning regime; they are not uniform high-dimensional
claims or optimality results. Source model approximation remains separate.
